\documentclass{article}

\title{Ch. 6.2 HW}
\author{Principles of Microeconomics}
\date{\today}

\begin{document}

\maketitle

\subsection*{Instructions}
These questions are Quick Quiz \#5-8 and Problems and Applications \#4-8 in Ch. 6. Please turn in your work to Gradescope by the end of class.

\subsection*{Quick Quiz}

\begin{enumerate}
\setcounter{enumi}{4}

\item A \$1 per unit tax levied on consumers of a good is equivalent to
	\begin{enumerate}
	
	\item a \$1 per unit tax levied on producers of the good.
	
	\item a \$1 per unit subsidy paid to producers of the good.
	
	\item a price floor that raises the good's price by \$1 per unit.
	
	\item a price ceiling that raises the good's price by \$1 per unit.
	
	\end{enumerate}
	
\item When a good is taxed, the burden falls mainly on consumers if
	\begin{enumerate}

	\item the tax is levied on consumers.
	
	\item the tax is levied on producers.
	
	\item supply is inelastic and demand is elastic.
	
	\item supply is elastic and demand is inelastic.
	
	\end{enumerate}
	
\item Which of the following increases the quantity supplied, decreases the quantity demanded, and increases the price that consumers pay?

	\begin{enumerate}
	
	\item the passage of a tax on a good
	
	\item the repeal of a tax on a good
	
	\item the imposition of a binding price floor
	
	\item the removal of a binding price floor
	
	\end{enumerate}	
	
\item Which of the following increases the quantity supplied, increases the quantity demanded, and decreases the price that consumers pay?

	\begin{enumerate}
	
	\item the passage of a tax on a good
	
	\item the repeal of a tax on a good
	
	\item the imposition of a binding price floor
	
	\item the removal of a binding price floor
	
	\end{enumerate}

\end{enumerate}

\subsection*{Problems and Applications}

\begin{enumerate}

\setcounter{enumi}{3}

\item Suppose the federal government requires beer drinkers to pay a \$2 tax on each case of beer purchased. (In fact, both the federal and state governments impose beer taxes of some sort.)

	\begin{enumerate}
	
	\item Draw a supply-and-demand diagram of the market for beer without the tax. Show the price paid by consumers, the price received by producers, and the quantity of beer sold. What is the difference between the price paid by consumers and the price received by producers?
	
	\item Now draw a supply-and-demand diagram for the beer market with the tax. Show the price paid by consumers, the price received by producers, and the quantity of beer sold. What is the difference between the price paid by consumers and the price received by producers? Has the quantity of beer sold increased or decreased?
	
	\end{enumerate}
	
\item A senator wants to raise tax revenue and make workers better off. A staff member proposes raising the payroll tax paid by firms and using part of the extra revenue to reduce the payroll tax paid by workers. Would this accomplish the senator's goal? Explain.

\item If the government places a \$500 tax on luxury cars, will the price paid by consumers rise by more than \$500, less than \$500, or exactly \$500? Explain.

\item Congress and the president decide that the United States should reduce air pollution by reducing its use of gasoline. They impose a \$0.50 tax on each gallon of gasoline sold.

	\begin{enumerate}
	
	\item Should they impose this tax on producers or consumers? Explain carefully using a supply-and-demand diagram.
	
	\item If the demand for gasoline were more elastic, would this tax be more effective or less effective in reducing the quantity of gasoline consumed? Explain with both words and a diagram.
	
	\item Are consumers of gasoline helped or hurt by this tax? Why?
	
	\item Are workers in the oil industry helped or hurt by this tax? Why?
	
	\end{enumerate}
	
\item A case study in this chapter discusses the federal minimum-wage law.

	\begin{enumerate}

	\item Suppose the minimum wage is above the equilibrium wage in the market for unskilled labor. Using a supply-and-demand diagram of the market for unskilled labor, show the market wage, the number of workers who are employed, and the number of workers who are unemployed. Also, show the total wage payments to unskilled workers.
	
	\item Now suppose the Secretary of Labor proposes an increase in the minimum wage. What effect would this increase have on employment? Does the change in employment depend on the elasticity of demand, the elasticity of supply, both elasticities, or neither?
	
	\item What effect would this increase in the minimum wage have on unemployment? Does the change in unemployment depend on the elasticity of demand, the elasticity of supply, both elasticities, or neither?
	
	\item If the demand for unskilled labor were inelastic, would the proposed increase in the minimum wage raise or lower total  wage payments to unskilled workers? Would your answer change if the demand for unskilled labor were elastic?
	
	\end{enumerate}
	
\end{enumerate}

\end{document}